\section{Biến đổi đơn giản và rút gọn biểu thức chứa căn thức bậc hai}

\begin{lesson}
    \subsection{Kiến thức trọng tâm}

    Không sử dụng MTCT, có thể so sánh được hai số $a=3\sqrt{2}$ và $b=2\sqrt{3}$ hay không?

    \answerline

    \subsubsection{Đưa thừa số ra ngoài dấu căn}

    \textbf{Cách đưa thừa số ra ngoài dấu căn}

    \textbf{HĐ1.} Tính và so sánh $\sqrt{(-3)^2\cdot25}$ với $|-3|\cdot\sqrt{25}$.

    \answerline

    \begin{keyconcept}
        Nếu $a$ là một số và $b$ là một số không âm thì
        \[
            \sqrt{a^2\cdot b}=|a|\sqrt{b}.
        \]
    \end{keyconcept}

    \textbf{Chú ý.} Phép biến đổi trên gọi là phép đưa thừa số ra ngoài dấu căn.

    \textbf{Ví dụ 1.} Viết nhân tử số của biểu thức dưới dấu căn thành tích các lũy thừa rồi đưa thừa số ra ngoài dấu căn:
    \begin{enumerate}[label=\alph*)]
        \item $\sqrt{45}$;
        \item $\sqrt{243a}$ $(a\ge0)$.
    \end{enumerate}

    \textbf{Giải}

    a) Ta có $45=3^2\cdot5$ nên
    \[
        \sqrt{45}=\sqrt{3^2\cdot5}=3\sqrt5;
    \]

    b) Ta có $243=3\cdot9^2$ nên
    \[
        \sqrt{243a}=\sqrt{9^2\cdot3a}=9\sqrt{3a}.
    \]

    \textbf{Luyện tập 1.} Đưa thừa số ra ngoài dấu căn:
    \begin{enumerate}[label=\alph*)]
        \item $\sqrt{12}$;
        \item $3\sqrt{27}$;
        \item $5\sqrt{48}$.
    \end{enumerate}

    \answerline
    \answerline

    \textbf{Chú ý.} Khi tính toán với những căn thức bậc hai mà biểu thức dưới dấu căn có mẫu, ta thường khử mẫu của biểu thức lấy căn (tức là biến đổi căn thức bậc hai đó thành một biểu thức mà trong căn thức không còn mẫu) như trong ví dụ sau:

    \textbf{Ví dụ 2.} Khử mẫu của biểu thức lấy căn $\sqrt{\dfrac47}$.

    \textbf{Giải}

    Nhân cả tử và mẫu của biểu thức lấy căn với số $7$ và đưa thừa số ra ngoài dấu căn, ta được
    \[
        \sqrt{\dfrac47}
        =\sqrt{\dfrac{4\cdot7}{7^2}}
        =\sqrt{\left(\dfrac27\right)^2\cdot7}
        =\dfrac{2\sqrt7}{7}.
    \]

    \textbf{Luyện tập 2.} Khử mẫu của biểu thức lấy căn $\sqrt{\dfrac35}$.

    \answerline

    \textbf{Tranh luận.} Vuông làm:
    \[
        \sqrt{(-2)^2\cdot5}=-2\sqrt5.
    \]
    Em có đồng ý với cách làm của Vuông không? Vì sao?

    \answerline

    \subsubsection{Đưa thừa số vào trong dấu căn}

    \textbf{Cách đưa một thừa số vào trong dấu căn}

    \textbf{HĐ2.} Tính và so sánh:
    \begin{enumerate}[label=\alph*)]
        \item $5\sqrt4$ với $\sqrt{5^2\cdot4}$;
        \item $-5\sqrt4$ với $-\sqrt{(-5)^2\cdot4}$.
    \end{enumerate}

    \answerline

    \begin{keyconcept}
        \begin{itemize}
            \item Nếu $a$ và $b$ là hai số không âm thì
            \[
                a\sqrt b=\sqrt{a^2b}.
            \]
            \item Nếu $a$ là số âm và $b$ là số không âm thì
            \[
                a\sqrt b=-\sqrt{a^2b}.
            \]
        \end{itemize}
    \end{keyconcept}

    \textbf{Chú ý.} Các phép biến đổi trên gọi là phép đưa thừa số vào trong dấu căn.

    \textbf{Ví dụ 3.} Đưa thừa số vào trong dấu căn:
    \begin{enumerate}[label=\alph*)]
        \item $5\sqrt2$;
        \item $-2\sqrt a$ $(a\ge0)$.
    \end{enumerate}

    \textbf{Giải}
    \begin{enumerate}[label=\alph*)]
        \item $5\sqrt2=\sqrt{5^2\cdot2}=\sqrt{50}$;
        \item $-2\sqrt a=-\sqrt{2^2\cdot a}=-\sqrt{4a}$.
    \end{enumerate}

    \textbf{Ví dụ 4.} Trả lời câu hỏi trong tình huống mở đầu.

    \textbf{Giải}

    Ta có
    \[
        3\sqrt2=\sqrt{3^2\cdot2}=\sqrt{18};
        \qquad
        2\sqrt3=\sqrt{2^2\cdot3}=\sqrt{12}.
    \]
    Vì $12<18$ nên $2\sqrt3<3\sqrt2$.

    \textbf{Luyện tập 3.} Đưa thừa số vào trong dấu căn:
    \begin{enumerate}[label=\alph*)]
        \item $3\sqrt5$;
        \item $-2\sqrt7$.
    \end{enumerate}

    \answerline

    \subsubsection{Trục căn thức ở mẫu}

    Tính toán với các biểu thức có chứa căn ở mẫu thường phức tạp và ta thường tìm cách trục các căn thức ở mẫu (tức là biến đổi biểu thức thành một biểu thức mới không chứa căn ở mẫu).

    \textbf{Cách trục căn thức ở mẫu}

    \textbf{HĐ3.} Nhân cả tử và mẫu của biểu thức $\dfrac{3a}{2\sqrt2}$ với $\sqrt2$ và viết biểu thức nhận được dưới dạng không có căn thức ở mẫu.

    \answerline

    \textbf{HĐ4.} Cho hai biểu thức $\dfrac{-2}{\sqrt3+1}$ và $\dfrac1{\sqrt3-\sqrt2}$. Hãy thực hiện các yêu cầu sau để viết các biểu thức đó dưới dạng không có căn thức ở mẫu:
    \begin{enumerate}[label=\alph*)]
        \item Xác định biểu thức liên hợp của mẫu.
        \item Nhân cả tử và mẫu với biểu thức liên hợp của mẫu.
        \item Sử dụng hằng đẳng thức hiệu hai bình phương để rút gọn mẫu của biểu thức nhận được.
    \end{enumerate}

    \answerline
    \answerline

    \begin{keyconcept}
        \begin{itemize}
            \item Với các biểu thức $A$, $B$ và $B>0$, ta có
            \[
                \dfrac{A}{\sqrt B}=\dfrac{A\sqrt B}{B}.
            \]
            \item Với các biểu thức $A$, $B$, $C$ mà $A\ge0$, $A\ne B^2$, ta có
            \[
                \dfrac{C}{\sqrt A+B}=\dfrac{C(\sqrt A-B)}{A-B^2};
                \qquad
                \dfrac{C}{\sqrt A-B}=\dfrac{C(\sqrt A+B)}{A-B^2}.
            \]
            \item Với các biểu thức $A$, $B$, $C$ mà $A\ge0$, $B\ge0$, $A\ne B$, ta có
            \[
                \dfrac{C}{\sqrt A+\sqrt B}
                =\dfrac{C(\sqrt A-\sqrt B)}{A-B};
                \qquad
                \dfrac{C}{\sqrt A-\sqrt B}
                =\dfrac{C(\sqrt A+\sqrt B)}{A-B}.
            \]
        \end{itemize}
    \end{keyconcept}

    \textbf{Ví dụ 5.} Trục căn thức ở mẫu của các biểu thức:
    \begin{enumerate}[label=\alph*)]
        \item $\dfrac2{3\sqrt5}$;
        \item $\dfrac a{3-2\sqrt2}$.
    \end{enumerate}

    \textbf{Giải}

    a) Nhân cả tử và mẫu của biểu thức đã cho với $\sqrt5$, ta được
    \[
        \dfrac2{3\sqrt5}
        =\dfrac{2\sqrt5}{3(\sqrt5)^2}
        =\dfrac{2\sqrt5}{15}.
    \]

    b) Biểu thức liên hợp của mẫu là $3+2\sqrt2$. Nhân cả tử và mẫu của biểu thức đã cho với $3+2\sqrt2$, ta được
    \[
        \dfrac a{3-2\sqrt2}
        =\dfrac{a(3+2\sqrt2)}{(3-2\sqrt2)(3+2\sqrt2)}
        =\dfrac{a(3+2\sqrt2)}{3^2-(2\sqrt2)^2}
        =a(3+2\sqrt2).
    \]

    \textbf{Luyện tập 4.} Trục căn thức ở mẫu của các biểu thức sau:
    \begin{enumerate}[label=\alph*)]
        \item $\dfrac{-5\sqrt{x^2+1}}{2\sqrt3}$;
        \item $\dfrac{a^2-2a}{\sqrt a+\sqrt2}$ $(a\ge0,\ a\ne2)$.
    \end{enumerate}

    \answerline
    \answerline

    \subsubsection{Rút gọn biểu thức chứa căn thức bậc hai}

    Khi rút gọn biểu thức có chứa căn thức bậc hai, ta cần phối hợp các phép tính (cộng, trừ, nhân, chia) và các phép biến đổi đã học (đưa thừa số ra ngoài hoặc vào trong dấu căn; khử mẫu của biểu thức lấy căn; trục căn thức ở mẫu).

    \textbf{Ví dụ 6.} Rút gọn biểu thức
    \[
        A=2\sqrt3-\sqrt{75}+\sqrt{(1-\sqrt3)^2}.
    \]

    \textbf{Giải}

    Đưa thừa số ra ngoài dấu căn, ta có:
    \[
        \sqrt{75}=\sqrt{3\cdot5^2}=5\sqrt3;
        \qquad
        \sqrt{(1-\sqrt3)^2}=|1-\sqrt3|=\sqrt3-1.
    \]
    Do đó
    \[
        A=2\sqrt3-5\sqrt3+\sqrt3-1=-1-2\sqrt3.
    \]

    \textbf{Ví dụ 7.}
    \begin{enumerate}[label=\alph*)]
        \item Trục căn thức ở mẫu của các biểu thức
        \[
            \dfrac{x^2-1}{\sqrt x-1};
            \qquad
            \dfrac{x^2-x}{\sqrt x+1}
        \]
        với $x>1$.
        \item Sử dụng kết quả câu a, rút gọn biểu thức
        \[
            P=\dfrac{x^2-1}{\sqrt x-1}-\dfrac{x^2-x}{\sqrt x+1}
        \]
        với $x>1$.
    \end{enumerate}

    \textbf{Giải}

    a) Ta có
    \[
        \dfrac{x^2-1}{\sqrt x-1}
        =\dfrac{(x^2-1)(\sqrt x+1)}{x-1}
        =(x+1)(\sqrt x+1)
        =x\sqrt x+\sqrt x+x+1,
    \]
    \[
        \dfrac{x^2-x}{\sqrt x+1}
        =\dfrac{x(x-1)(\sqrt x-1)}{x-1}
        =x(\sqrt x-1)
        =x\sqrt x-x.
    \]

    b) Sử dụng kết quả câu a ta có
    \[
        P=x\sqrt x+\sqrt x+x+1-(x\sqrt x-x)
        =2x+\sqrt x+1.
    \]

    \textbf{Luyện tập 5.} Rút gọn biểu thức sau:
    \[
        \left(
            \dfrac{\sqrt{22}-\sqrt{11}}{1-\sqrt2}
            +\dfrac{\sqrt{21}-\sqrt7}{1-\sqrt3}
        \right)(\sqrt7-\sqrt{11}).
    \]

    \answerline
    \answerline

    \textbf{Vận dụng.} Trong thuyết tương đối, khối lượng $m$ (kg) của một vật khi chuyển động với tốc độ $v$ (m/s) được cho bởi công thức
    \[
        m=\dfrac{m_0}{\sqrt{1-\dfrac{v^2}{c^2}}},
    \]
    trong đó $m_0$ (kg) là khối lượng của vật khi đứng yên, $c$ (m/s) là tốc độ của ánh sáng trong chân không (Theo sách Vật lí đại cương, NXB Giáo dục Việt Nam, 2016).
    \begin{enumerate}[label=\alph*)]
        \item Viết lại công thức tính khối lượng $m$ dưới dạng không có căn thức ở mẫu.
        \item Tính khối lượng $m$ theo $m_0$ (làm tròn đến chữ số thập phân thứ ba) khi vật chuyển động với vận tốc $v=\dfrac1{10}c$.
    \end{enumerate}

    \answerline
    \answerline
\end{lesson}

\begin{exercises}
    \subsection{Bài tập}

    \begin{questions}[resume=SGK]
        \item Đưa thừa số ra ngoài dấu căn:
        \begin{questions}
            \item $\sqrt{52}$;
            \item $\sqrt{27a}$ $(a\ge0)$;
            \item $\sqrt{50\sqrt2+100}$;
            \item $\sqrt{9\sqrt5-18}$.
        \end{questions}

        \answerline
        \answerline

        \item Đưa thừa số vào trong dấu căn:
        \begin{questions}
            \item $4\sqrt3$;
            \item $-2\sqrt7$;
            \item $4\sqrt{\dfrac{15}{2}}$;
            \item $-5\sqrt{\dfrac{16}{5}}$.
        \end{questions}

        \answerline
        \answerline

        \item Khử mẫu trong dấu căn:
        \begin{questions}
            \item $2a\sqrt{\dfrac35}$;
            \item $-3x\sqrt{\dfrac5x}$ $(x>0)$;
            \item $-\sqrt{\dfrac{3a}{b}}$ $(a\ge0,\ b>0)$.
        \end{questions}

        \answerline
        \answerline

        \item Trục căn thức ở mẫu:
        \begin{questions}
            \item $\dfrac{4+3\sqrt5}{\sqrt5}$;
            \item $\dfrac1{\sqrt5-2}$;
            \item $\dfrac{3+\sqrt3}{1-\sqrt3}$;
            \item $\dfrac{\sqrt2}{\sqrt3+\sqrt2}$.
        \end{questions}

        \answerline
        \answerline

        \item Rút gọn các biểu thức sau:
        \begin{questions}
            \item $2\sqrt{\dfrac23}-4\sqrt{\dfrac32}$;
            \item $\dfrac{5\sqrt{48}-3\sqrt{27}+2\sqrt{12}}{\sqrt3}$;
            \item $\dfrac1{3+2\sqrt2}+\dfrac{4\sqrt2-4}{2-\sqrt2}$.
        \end{questions}

        \answerline
        \answerline

        \item Rút gọn biểu thức
        \[
            A=\sqrt x\left(\dfrac1{\sqrt x+3}-\dfrac1{3-\sqrt x}\right)
            \qquad (x\ge0,\ x\ne9).
        \]

        \answerline
        \answerline
    \end{questions}
\end{exercises}

\begin{practice}
    \subsection{Luyện tập}

    \begin{questions}[resume=SBT]
        \item Tính giá trị biểu thức
        \[
            P=(\sqrt{20}+2\sqrt{45}-3\sqrt{80})^2.
        \]

        \answerline
        \answerline

        \item Sắp xếp các số sau theo thứ tự tăng dần:
        \[
            9\sqrt2;\quad 8\sqrt3;\quad 5\sqrt6;\quad 4\sqrt7.
        \]

        \answerline

        \item Thực hiện phép tính
        \[
            \left(\dfrac1{\sqrt8+\sqrt7}+\sqrt{175}-2\sqrt2\right)^2.
        \]

        \answerline
        \answerline

        \item Thực hiện phép tính
        \[
            \sqrt{\dfrac{3-2\sqrt2}{17-12\sqrt2}}
            -\sqrt{\dfrac{3+2\sqrt2}{17+12\sqrt2}}.
        \]

        \answerline
        \answerline

        \item Không sử dụng MTCT, chứng minh rằng biểu thức sau có giá trị là một số nguyên:
        \[
            P=
            \left(
                \dfrac{\sqrt5+1}{1+\sqrt5+\sqrt3}
                +\dfrac{\sqrt5-1}{1+\sqrt3-\sqrt5}
            \right)
            \left(\sqrt3-\dfrac4{\sqrt3}+2\right)
            \sqrt{0,2}.
        \]

        \answerline
        \answerline
        \answerline

        \item
        \begin{questions}
            \item Trục căn ở mẫu của biểu thức
            \[
                \dfrac{3+\sqrt2}{2\sqrt2-1}.
            \]
            \item Tính giá trị biểu thức
            \[
                P=x(x^4-6x^2+1)
            \]
            tại
            \[
                x=\dfrac{3+\sqrt2}{2\sqrt2-1}.
            \]
        \end{questions}

        \answerline
        \answerline
    \end{questions}
\end{practice}
